% Figura redibujada a partir de "Imagenes/Diagrama de bandas VGS g 0 VDS g 0.PNG" (Tema 5).
% Coordenadas en píxeles del PNG original (820 x 691 px), escaladas para que
% ocupe lo mismo que el PNG en el documento: 246.9 pt x 208.1 pt.
\documentclass[12pt,border=0pt]{standalone}
\usepackage[spanish,es-tabla]{babel}
\usepackage[utf8]{inputenc}
\usepackage{amsmath}
\usepackage{xcolor}
\usepackage{tikz}
\usepackage{pgfplots}
\usepackage{circuitikz}
\usetikzlibrary{babel, circuits.ee}

\definecolor{applegreen}{rgb}{0.55, 0.71, 0.0}
\definecolor{airforceblue}{rgb}{0.36, 0.54, 0.66}
\definecolor{amber}{rgb}{1.0, 0.49, 0.0}

\newcommand{\diff}[3][1]{%
	\IfBlankTF{#1}{\frac{\mathrm{d}{#2}}{\mathrm{d}{#3}}}%
	{\frac{\mathrm{d}^{#1} {#2}}{\mathrm{d}{#3}^{#1}}}%
}%
\newcommand{\pdiff}[3][1]{%
	\IfBlankTF{#1}{\frac{\partial {#2}}{\partial {#3}}}%
	{\frac{\partial^{#1} {#2}}{\partial {#3}^{#1}}}%
}%

\begin{document}
\begin{tikzpicture}[x=0.30110pt, y=-0.30110pt, font=\footnotesize,
    banda/.style={ultra thick, line join=round},
    fermi/.style={blue, thick, dashed, line join=round},
    intrinseco/.style={thick, dotted},
    eje/.style={}]
    \useasboundingbox (0,0) rectangle (820,691);
    % ---- Ejes ----
    \draw[-latex] (182,690) -- (182,78);
    \draw[-latex] (0,663) -- (785,663);
    \draw (477,95) -- (477,690);
    \draw[-latex] (692,640) -- (752,588);
    \node at (125,48) {Energía};
    \node at (225,636) {$y=0$};
    \node at (430,636) {$y=L$};
    \node at (802,650) {$y$};
    \node at (776,595) {$x$};
    % ---- E_C ----
    \draw[banda] (20,220) -- (150,220) to[out=0, in=180] (210,185)
        plot[smooth, tension=0.6] coordinates {(210,185) (245,187) (300,205) (350,232)
            (400,282) (440,340) (470,400) (500,447)}
        -- (640,447) -- (755,345) -- (645,345);
    % Paredes de la región P (volumen)
    \draw[banda] plot[smooth, tension=0.6] coordinates
        {(195,183) (210,150) (233,98) (262,60) (295,32) (315,13)};
    \draw[banda] (315,13) -- (553,13);
    \draw[banda] plot[smooth, tension=0.6] coordinates
        {(467,392) (474,320) (488,220) (508,140) (530,70) (553,13)};
    \draw[banda] plot[smooth, tension=0.6] coordinates
        {(553,13) (566,22) (580,62) (593,150) (604,240) (614,300) (626,332) (645,345)};
    \draw (530,418) -- (615,348);
    % ---- Cuasiniveles de Fermi ----
    % E_FN
    \draw[fermi] (20,245) -- (230,245)
        plot[smooth, tension=0.6] coordinates {(230,245) (290,272) (345,322) (388,380)
            (420,435) (458,468) (500,480)}
        -- (640,480) -- (755,375);
    % E_FP
    \draw[fermi] (230,245) -- (510,245) to[out=0, in=90] (522,262) -- (522,420)
        to[out=270, in=110] (540,470);
    \draw[fermi] (372,245) -- (540,180) -- (598,133);
    \draw[fermi] (22,243) -- (150,133) -- (598,133) to[out=0, in=90] (633,155)
        -- (633,330) to[out=270, in=180] (648,345);
    % ---- E_i ----
    \draw[intrinseco] (20,315) -- (145,315) to[out=0, in=180] (205,282)
        plot[smooth, tension=0.6] coordinates {(205,282) (235,280) (285,287) (335,307)
            (372,333) (405,375) (440,432) (470,490) (500,548)}
        -- (640,548) -- (755,438) -- (685,438);
    % ---- E_V ----
    \draw[banda] (20,408) -- (140,408) to[out=0, in=180] (205,375)
        plot[smooth, tension=0.6] coordinates {(205,375) (235,373) (285,382) (335,402)
            (382,433) (422,483) (460,555) (488,610) (505,640)}
        -- (640,640) -- (755,538) -- (650,538);
    % ---- Etiquetas ----
    \node at (395,42) {$P$};
    \node at (80,165) {$N^+$};
    \node at (625,395) {$N^+$};
    \node[blue] at (418,186) {$E_{\mathrm{F}P}$};
    \node[blue] at (368,382) {$E_{\mathrm{F}N}$};
    \node at (786,325) {$E_\mathrm{C}$};
    \node[blue] at (785,372) {$E_{\mathrm{F}N}$};
    \node at (782,440) {$E_\mathrm{Fi}$};
    \node at (786,537) {$E_\mathrm{V}$};
    \node at (65,478) {Fuente};
    \node at (298,488) {Sustrato};
    \node at (572,580) {Drenador};
\end{tikzpicture}
\end{document}
